\section{Solving the Master Program at Scale: Shape Aggregation}
\label{app:shape_agg}

The comprehensive V3 program (\cref{app:master_v3}) is large: built naively, with
one copy of every tiling, coupling, and roofline variable per operation, the
DeepSeek-V3 instance reaches ${\sim}8.7$M rows and does not solve within a
practical iteration budget. This appendix describes the scalability path we adopt,
adapted from the resource-allocation solver of \citet{cvxcluster2026}, which solves
granular cluster placement $100$--$1000\times$ faster by aggregating identical
entities into \emph{shapes} rather than reasoning about each entity individually.
The same principle collapses our program by roughly two orders of magnitude.

\paragraph{Servers, shapes, and multiplicities.}
In the cluster-placement setting of \citet{cvxcluster2026}, a \emph{server} is one
physical placement target, and a \emph{shape} is an equivalence class of servers
with identical resource profiles (the same CPU/memory/GPU configuration). A cluster
with thousands of servers has only tens of distinct shapes, so one aggregated
variable per shape, weighted by its \emph{multiplicity} (how many servers it
stands for), replaces thousands of per-server variables. The linear program then
scales with the number of \emph{shapes}, not the number of \emph{servers}.

The analogy to our program is exact once the entities are named correctly. Our
``server'' is an \emph{operation instance}: a node at a specific position in the
network's execution order. Our ``shape'' is an equivalence class of operations with
\emph{identical dimensions} $(C,M,R,S,P,Q,N)$ and strides. Two operations of the
same shape have byte-for-byte identical Timeloop Pareto data and roofline
coefficients (\cref{sec:v3_levers} confirmed this: our four target LLMs comprise
$720$--$903$ operations each but only $29$ \emph{unique} shapes in total; Llama-3
70B's $720$ operations collapse to just $6$). Every shape-dependent block of the
program (tile selection $y$, the McCormick auxiliaries $w,\zeta$, the per-tile
compute and memory coefficients $\Tmin,\Tmax,t_i^k,d^{\ast}$, and the perspective
epigraph $\tau$) is therefore identical across all operations of a shape. Building
one aggregated copy per shape $s$ with multiplicity $m_s$ (the count of operations
of that shape), and replacing $\sum_i T_i$ by $\sum_s m_s\, T_s$, shrinks these
blocks from $O(N)$ to $O(S)$ with $S \ll N$ (here $S\approx 6$ versus $N=720$), the
${\sim}100\times$ reduction that makes the single-solve program tractable.

\paragraph{Decoupling the shape dimension from the time dimension.}
There is one difference between our program and cluster placement that must be
handled with care, and it is the crux of the aggregation. The placement problem of
\citet{cvxcluster2026} is solved as a sequence of \emph{static snapshots}: the
program at any one instant has no notion of time, so identical servers are fully
interchangeable within it and aggregation is lossless. (Their scheduler is
online and re-solves as the workload mix changes; it is each round's program,
not the system, that is time-free.) Our
program has a \emph{time-indexed schedule}: the residency variables $p_{e,t}$ and
the per-timestep capacity rows track which tensors are pinned, evicted, or
recomputed at each step, and two operations of the same shape sitting at
\emph{different} positions in the execution order have different residency states.
Collapsing the time axis by shape would therefore be incorrect.

We resolve this by \emph{factoring the program along its two axes} and aggregating
only the axis that permits it. The \textbf{shape-invariant block}, tiling, the
compute and memory rooflines, the McCormick coupling, and the entire hardware-sizing
layer (the reciprocals $v,u$, the caps, and every decision of \cref{sec:v3_levers})
,  depends on an operation only through its shape, and is aggregated per unique
shape. This block is where the $8.7$M rows lived: they were overwhelmingly redundant
copies of the same handful of shapes, and aggregating them is the entire size win.
The \textbf{time-indexed residency block}, $p_{e,t}$, rematerialization, and the
capacity rows, stays per-instance, because it encodes \emph{position}, not shape;
it is a comparatively small, banded structure, and it can be compressed further on
its own terms (a steady-state residency pattern per shape, or a coarsened time grid)
without touching the shape-aggregated part. Concretely: ``decoupling server from
time'' means writing the program as (shape-dependent, aggregatable) rows times
(time-dependent, per-instance) rows, and applying shape aggregation to the first
factor only. The hardware variables $v,u,\mathrm{MACs},\mathrm{BW},C_{\mathrm{GM}}$
are already global (one set for the whole chip), so the entire co-design layer is
shape- and time-invariant and aggregates trivially.

\paragraph{Dual-guided rounding.}
\citet{cvxcluster2026} also extract the LP's dual variables, dual variables on
each resource, and use them to rank and greedily round the relaxed assignments,
which is tighter than naive rounding. We already round the relaxed tile selection by
$\arg\max_c y_{i,c}$ and re-solve (\cref{sec:v1}); the same dual-guided ranking
sharpens the rounding of the \emph{new} integer knobs the comprehensive program adds
,  the expert-pin $g_e$, the precision one-hots $q_{\rho}$, and the array-aspect
one-hot $\alpha_a$, by pinning first the knobs whose reduced cost most improves
the certified bound.

\paragraph{An independent validation of the paradigm.}
Beyond the scalability trick, \citet{cvxcluster2026} is a useful external data point:
a contemporaneous system that, for a completely different problem (datacenter task
placement), independently arrives at CHEETAH's exact recipe, relax a hard discrete
program to a convex one, read the dual variables as interpretable dual signals that
feed an outer decision loop, and round greedily against those duals. That the same
LP-relaxation-plus-dual-feedback structure is effective across such different
domains is evidence that the paradigm at the core of CHEETAH is sound and general,
rather than an artifact of accelerator co-design.
